\hypersetup{pageanchor=false}
\vspace*{0cm}%
\FTMlayoutPutAtNorthWest{2.5cm}{1.5cm}{\FTMlogoTercetChair{1.05cm}}%
\FTMlayoutPutAtNorthEast{2.4cm}{2cm}{\FTMlogoTUM{1cm}{TUMBlue}}%

{\fontsize{24}{24} \selectfont \FTMlangProjectDescription}
\vspace{1cm}

{\fontsize{16}{18} \selectfont Titel der Arbeit}
\vspace{0,5cm}

Hier steht ein zum Thema hinführender Text.

In einer konstruktiven (oder theoretischen) \textcolor{red}{Bachelor/-Semester/-Masterarbeit} sollen (...)

Folgende Punkte sind durch \textcolor{red}{Herrn/Frau Martin Musterstudent} zu bearbeiten:

\begin{itemize}
	\item Punkt 1
	\item Punkt 2
\end{itemize}

Die Ausarbeitung soll die einzelnen Arbeitsschritte in übersichtlicher Form dokumentieren. Der Kandidat/Die Kandidatin verpflichtet sich, die \textcolor{red}{Bachelor/-Semester/-Masterarbeit} selbständig durchzuführen und die von ihm verwendeten wissenschaftlichen Hilfsmittel anzugeben. 

Die eingereichte Arbeit verbleibt als Prüfungsunterlage im Eigentum des Lehrstuhls.
\vspace{1cm}

\begin{longtable}{p{0.5\linewidth}p{0.44\linewidth}}
	\hspace{-0.2cm}\normalsize \FTMlangAnnouncementDate: \FTMutilsDate{1}{4}{2021} & \hspace{1.0cm}\normalsize \FTMlangSubmissionDate: \FTMutilsDate{1}{10}{2021}\\
\end{longtable}
\vspace{0.5cm}
\begin{longtable}{p{0.5\linewidth}p{0.44\linewidth}} 
	\normalsize
	\hspace{-0.2cm}\rule{5.8cm}{1pt} & \hspace{1.0cm}\rule{5.8cm}{1pt}\\
	\\
	\hspace{-0.2cm}\normalsize \FTMnamesProfLienkamp & \hspace{1.0cm}\normalsize Max Musterbetreuer, M.Sc.\\
\end{longtable}

%% To directly include a PDF-File, delete the text above and use:
% \includepdf[pages={1}]{./source/Aufgabenstellung_BASAMA.pdf}

\cleardoublepage
\hypersetup{pageanchor=true}